% SEGMENT OLP-0399-S01
% Part: many-valued-logic
% Chapter: infinite-valued-logics
% Section: introduction

\documentclass[../../../include/open-logic-section]{subfiles}

\begin{document}

\olfileid{mvl}{inf}{int}

\olsection{அறிமுகம்}

ஓர் அணியிலுள்ள மெய்மதிப்புகளின் எண்ணிக்கை முடிவுறுவதாக இருக்க வேண்டியதில்லை.
முடிவிலாப் பல மெய்மதிப்புகளைக் கொண்ட ஒரு கணத்திற்கான இயல்பான தேர்வு,
$0$ க்கும்~$1$ க்கும் இடையிலுள்ள விகிதமுறு எண்களின் கணம்
$V_\infty = [0,1] \cap \Rat$; அதாவது,
% SOURCE CORRECTION TA-MVL-016: use n<m so V_m has exactly m values in [0,1].
\begin{align*}
    V_\infty & = \Setabs{\frac{n}{m}}{n,m \in \Nat \text{ மற்றும் } n\le m}.
\intertext{இந்த முடிவிலா மெய்மதிப்புக் கணத்தைக் கருதும்போது, பின்வரும்
உட்கணங்களையும் கருதுவது பல நேரங்களில் பயனுள்ளது:}
V_m & = \Setabs{\frac{n}{m-1}}{n \in \Nat \text{ மற்றும் } n<m}
\intertext{எடுத்துக்காட்டாக, $V_5$ என்பது சீரான இடைவெளியில் அமைந்த $5$
மெய்மதிப்புகளைக் கொண்ட கணம்:}
V_5 & = \{0, \frac{1}{4}, \frac{1}{2}, \frac{3}{4}, 1\}.
\end{align*}
இந்த மெய்மதிப்புக் கணங்களை அடிப்படையாகக் கொண்ட தருக்கங்களில், வழக்கமாக
$1$ மட்டுமே நியமிக்கப்படுகிறது; அதாவது, $V^+ = \{1\}$. வேறு சொற்களில்,
$1$ முழுமையான மெய்மையின் இடத்தையும், $0$ முழுமையான பொய்மையின் இடத்தையும்
வகிக்குமாறு செய்கிறோம்; ஆனால் !!{formula}s ஆனவை~$V$ இல் உள்ள எந்த
இடைநிலை மதிப்பையும் பெறலாம்.

எனினும், $0$ க்கும்~$1$ க்கும் இடையிலுள்ள எல்லா \emph{மெய்யெண்களையும்}
கொண்ட $V_{[0,1]} = [0,1]$ என்ற கணத்தையோ, $[0,1]$ இன் வேறு முடிவிலா
உட்கணங்களையோ கருதலாம். இந்த மெய்மதிப்புக் கணத்தைக் கொண்ட தருக்கங்கள்
பெரும்பாலும் \emph{மங்கல் தருக்கங்கள்} எனப்படுகின்றன.


\end{document}
